\documentclass[11pt,a4paper]{article}
\usepackage[utf8]{inputenc}
\usepackage[T1]{fontenc}
\usepackage{lmodern}
\usepackage[margin=2.1cm]{geometry}
\usepackage{booktabs}
\usepackage{amsmath}
\usepackage{graphicx}
\usepackage{xcolor}
\usepackage{hyperref}
\usepackage{microtype}
\usepackage{caption}
\usepackage{float}
\hypersetup{colorlinks=true,linkcolor=blue!55!black,urlcolor=blue!55!black,
  pdftitle={Phase 2: Leakage-Safe Spatial and Comparable-Sale Features for Paris Apartment Valuation},
  pdfauthor={Badr Ettousy}}
\graphicspath{{../../visuals/phase2/}}
\setlength{\parskip}{0.35em}
\setlength{\parindent}{0pt}

\title{Phase 2: Leakage-Safe Spatial and Comparable-Sale Features\\
for Paris Apartment Valuation}
\author{Badr Ettousy}
\date{August 2026}

\begin{document}
\maketitle

\begin{abstract}
This addendum evaluates a second-stage data-engineering pipeline for the Paris
DVF apartment valuation benchmark. A typed Gold table of 143,009 transactions
adds Lambert-93 distances, H3 cells, local transaction density, and historical
comparable-sale statistics. Every comparable is at least 90 days older than the
transaction being valued. Model and blend selection use 2024 only; 28,330 sales
from 2025 form the final chronological test. A transparent correction combining
75\% of the Phase 1 LightGBM estimate with 25\% of a trailing local median
reduces test MAE from EUR 96,981 to EUR 96,428. The paired bootstrap interval
for the reduction is EUR 333--743. Full enriched LightGBM and XGBoost models do
not materially improve the Phase 1 reference. The result supports conservative
use of recent local evidence and argues against complexity without new
building-level information.
\end{abstract}

\section{Objective and contribution}
Phase 1 established a chronological six-model benchmark. Phase 2 asks whether
better spatial data engineering and strictly historical comparable sales add
out-of-year value. The contribution is an auditable feature contract, not a new
learning algorithm. Source-file hashes, cleaning counts, feature ranges, and
coverage are recorded in a machine-readable quality report.

\section{Data-engineering design}
Five annual Paris DVF files covering 2021--2025 are cleaned at mutation level
using the fixed Phase 1 cohort definition. WGS84 coordinates are retained for
storage and transformed to Lambert-93 for metric distance calculations. H3
cells at resolutions 8, 9, and 10 preserve multiple spatial scales. Address,
street, commune, and parcel identifiers are retained in the Gold table, while
the highest-cardinality categories are not supplied directly to the models.

For each sale at date $t$, eligible historical observations satisfy
\[
  90 \leq t-t_i \leq 730 \text{ days}, \qquad d(i,t) \leq 1000 \text{ m}.
\]
Counts and median prices per square metre are computed for several nested
distance/time windows. Additional features include dispersion, distance/time
weighted comparable values, size/room-similar comparable values, nearest prior
sale distance, median comparable age, and short-versus-long local momentum.
The 90-day lag approximates information availability and prevents same-period
or future transactions from leaking into a valuation.

\begin{figure}[H]
\centering
\includegraphics[width=0.98\textwidth]{comparable_coverage.png}
\caption{Coverage and density of eligible historical evidence. Lower 2021
coverage is expected because no pre-2021 source file is available.}
\end{figure}

\section{Evaluation protocol}
Development training uses 90,275 sales from 2021--2023. The 24,404 sales from
2024 select model configurations, target formulations, comparable definitions,
and blend weights. Selected learners are refitted on 114,679 sales from
2021--2024 and tested once on 28,330 sales from 2025. The primary metric is MAE;
secondary metrics are median absolute percentage error (MdAPE), $R^2$, and the
share of estimates within 20\% of the transaction value.

LightGBM and XGBoost use full spatial/comparable features with log total-price,
log-price-per-square-metre, and local-residual targets. Spatial-only and
comparable-only ablations test the contribution of feature groups. A six-model
blend and four transparent Phase 1/comparable corrections are selected using
validation MAE only.

\begin{figure}[H]
\centering
\includegraphics[width=0.98\textwidth]{comparable_weight_selection.png}
\caption{Validation-only selection of the comparable feature and blend weight.
The selected correction assigns 75\% weight to Phase 1 and 25\% to the
500-metre, trailing-365-day median.}
\end{figure}

\section{Results}
\begin{table}[H]
\centering
\caption{Selected 2025 test results. Lower MAE and MdAPE are better.}
\begin{tabular}{lrrr}
\toprule
Method & MAE (EUR) & MdAPE & Within 20\% \\
\midrule
Phase 1 + comparable correction & \textbf{96,428} & \textbf{12.22\%} & \textbf{71.69\%} \\
XGBoost, total-price target & 96,914 & 12.43\% & 71.18\% \\
Phase 1 LightGBM & 96,981 & 12.33\% & 71.41\% \\
XGBoost, EUR/m$^2$ target & 96,996 & 12.61\% & 70.73\% \\
Six-model validation blend & 97,183 & 12.52\% & 70.89\% \\
XGBoost, residual target & 97,507 & 12.60\% & 70.69\% \\
LightGBM, spatial-only & 98,232 & 12.72\% & 70.31\% \\
LightGBM, comparable-only & 98,963 & 12.73\% & 70.81\% \\
LightGBM, full total-price & 99,144 & 12.73\% & 69.89\% \\
\bottomrule
\end{tabular}
\end{table}

\begin{figure}[H]
\centering
\includegraphics[width=\textwidth]{phase2_accuracy.png}
\caption{Chronological test accuracy. Error bars are 95\% bootstrap intervals
for each MAE.}
\end{figure}

The selected correction improves MAE by EUR 552 relative to Phase 1. A paired
bootstrap over identical test transactions gives a 95\% interval of EUR
333--743 for this reduction. MdAPE improves from 12.33\% to 12.22\%, and the
within-20\% rate increases from 71.41\% to 71.69\%. $R^2$ decreases slightly
from 0.8619 to 0.8601, so the evidence is metric-dependent. The effect is
statistically consistent on this cohort but modest in practical magnitude.

\begin{figure}[H]
\centering
\includegraphics[width=\textwidth]{correction_diagnostics.png}
\caption{Prediction, error-distribution, temporal-stability, and paired-error
diagnostics for the selected correction.}
\end{figure}

\section{Interpretation}
Adding all engineered variables to a flexible learner did not outperform the
simpler Phase 1 model. The feature-gain diagnostic confirms that surface remains
dominant, although XGBoost uses the local median and weighted comparable
features. Importance values describe model usage, not causal effects.

\begin{figure}[H]
\centering
\includegraphics[width=\textwidth]{phase2_feature_importance.png}
\caption{Gain-based feature importance for the full total-price models.}
\end{figure}

The transparent correction succeeds by shrinking the base estimate only partly
toward recent local evidence. It is easier to inspect and reproduce than the
six-model blend. The outcome also shows why engineering effort should be judged
by chronological holdout performance rather than feature count.

\section{Limitations and next step}
DVF lacks condition, floor, lift, view, noise, renovation quality, detailed
building state, and energy performance. Because 2025 had already been inspected
during Phase 1, future claims require genuinely prospective confirmation on a
later cohort. The bootstrap quantifies sampling variation within 2025, not
market-regime uncertainty. The estimates are statistical indications, not
certified appraisals or investment guarantees.

The recommended Phase 3 is reliable external linkage: BAN-normalized addresses,
BDNB building characteristics, DPE ratings, cadastral age, and distances to
transit and amenities. These additions address missing information directly and
should be assessed with both chronological and spatial holdouts before exploring
attention-based architectures.

\section{Reproducibility}
The complete build, benchmark, visualization, and single-property inference
commands are documented in \texttt{reports/phase2/PHASE2\_REPORT.md}. Numerical
results are stored as CSV and JSON, and the Gold table is Parquet with its source
hashes and feature contract recorded separately.

\end{document}
